%==============================
%	Importation des packages
%==============================
% --- Encodage & langue ---
\usepackage[utf8]{inputenc}
\usepackage[T1]{fontenc}
\usepackage[french]{babel}

% --- Maths ---
\usepackage{amsmath, amssymb, amsthm, amsfonts}
\usepackage{bm}

% --- Mise en page ---
\usepackage[margin=1in]{geometry}
\usepackage{lmodern}
\usepackage[colorlinks=true, linkcolor=DAcolor, urlcolor=DAcolor]{hyperref}
\usepackage{graphicx}
\usepackage{todonotes}
\usepackage{marginnote}
\usepackage{fancyhdr}
\usepackage{titlesec}

\usepackage{setspace}


\usepackage{tcolorbox}
\tcbuselibrary{theorems, skins}
\usepackage{xcolor}

\tcbuselibrary{skins, breakable}

\usepackage{hyperref}

%===================================
%	Implémentation des lectures
%===================================
\setuptodonotes{noline, color=DAcolor!20, size=\footnotesize}
\makeatletter
\@todonotes@SetTodoListName{Cours}
\makeatother

\newcounter{lecture}


%=================
%	Commandes
%=================
% Ensembles
\newcommand{\N}{\mathbb{N}}
\newcommand{\Z}{\mathbb{Z}}
\newcommand{\Q}{\mathbb{Q}}
\newcommand{\R}{\mathbb{R}}
\newcommand{\C}{\mathbb{C}}
\newcommand{\K}{\mathbb{K}}

\DeclareMathOperator{\CL}{CL}
\DeclareMathOperator{\Hom}{Hom}
\DeclareMathOperator{\End}{End}
\DeclareMathOperator{\Vect}{Vect}
\DeclareMathOperator{\Span}{span}

\newcommand{\lecture}{\stepcounter{lecture}\todo{Lecture \thelecture}}

\newcommand{\et}{\quad \text{et} \quad}

\newcommand{\prf}[1]{
	\begin{proof}
		#1
	\end{proof}
	\vspace{0.5cm}
}



%========================================
%	Implémentation des environnements
%========================================
\newtheoremstyle{coloredtitle}
{}{}                          % espace avant/après
{\itshape}                    % corps en italique (comme plain)
{}                            % pas d'indentation
{\bfseries}                   % titre en gras (comme plain)
{.}                           % ponctuation après le titre
{ }                           % espace après
{\textcolor{DAcolor}{\bfseries\thmname{#1}\thmnumber{ #2}\thmnote{ (#3)}}}

\newtheoremstyle{coloredtitlenoint}
{}{}                          % espace avant/après
{}                    % corps en italique (comme plain)
{}                            % pas d'indentation
{\bfseries}                   % titre en gras (comme plain)
{.}                           % ponctuation après le titre
{ }                           % espace après
{\textcolor{DAcolor}{\bfseries\thmname{#1}\thmnumber{ #2}\thmnote{ (#3)}}}

\theoremstyle{coloredtitle}
\newtheorem{thm}{Théorème}[chapter]
\newtheorem{lem}{Lemme}[chapter]
\newtheorem{prop}{Proposition}[chapter]

\theoremstyle{coloredtitlenoint}
\newtheorem{definition}{Définition}[chapter]
\newtheorem{exemple}{Exemple}[chapter]

\setcounter{secnumdepth}{3}

% Théorème — encadré complet (peu fréquent, justifié)
\definecolor{DAcolor}{RGB}{50,0,100}    % théorème
\definecolor{DAcolordarker}{RGB}{25,0,50}
%\definecolor{DAcolor}{RGB}{0,0,102}    % théorème
%\definecolor{DAcolordarker}{RGB}{0,0,50}


% Définition — trait gauche uniquement (supporte les successions)



%=====================================
%	Mise en page et structure
%=====================================
\pagestyle{fancy}
\fancyhf{}  % efface tout

% Numéro de page en coin extérieur bas
\fancyfoot[LE]{\thepage}  % Left Even  → page de gauche
\fancyfoot[RO]{\thepage}  % Right Odd  → page de droite

% Optionnel — nom du cours en coin intérieur haut
\fancyhead[RE, LO]{\small\textit{\matiere}}
\renewcommand{\headrulewidth}{0pt}

%inter-ligne
\setstretch{1.1}

%configuration de l'indentation
\setlength{\parindent}{0pt}

%configuration de la numérotaion des équations
\numberwithin{equation}{section}

%configuration des sections et sous-sections
\titleformat{\section}[block]   % format "block"
{\centering\large\bfseries\color{DAcolordarker}}     % centrage + taille + gras
{\thesection}                 % numéro de la sous-section
{1em}                            % espace entre numéro et texte
{}  

\titleformat{\chapter}    % section = "Chapitre"
[block]                 % format en bloc
{\centering\large\color{DAcolordarker}}      % style : centré, grand (pas en gras)
{\MakeUppercase{Chapitre \thechapter}} % numéro en majuscules
{0pt}                   % pas d’espace entre numéro et texte
{\vspace{1ex}\\\bfseries\huge}  % saut de ligne + texte en énorme

\titlespacing*{\subsection}{0pt}{3.5ex plus 1ex minus .2ex}{3ex} 

\titleformat{\subsection}[block]
{\large\bfseries\color{DAcolordarker}}
{\thesubsection}
{1em}
{}

